\documentclass[10pt,a4paper]{amsart}
\usepackage[margin=19mm]{geometry}
\usepackage{amsmath,amssymb,mathtools}
\usepackage[scr=rsfs,cal=euler]{mathalfa}
\usepackage{xfrac}
\usepackage[unicode,hidelinks,bookmarks=false]{hyperref}
\newcommand{\ie}{i.e.\ }
\newcommand{\cf}{cf.\ }
\newcommand{\resp}{respectively\ }
\newcommand{\Subsection}{\subsection}
\providecommand{\nobreakdash}{\nobreak\hbox{-}}
\setlength{\emergencystretch}{2em}
\multlinegap0pt
\usepackage{bm}
\usepackage{bbm}
\usepackage{dsfont}
\usepackage{xspace}
\usepackage{cancel}
\usepackage{mathtools}
\usepackage{amsthm}
\makeatletter
\@ifundefined{theorem}{\newtheorem{theorem}{Theorem}}{}
\makeatother
\title[On the Transfer of Completeness and Projection Properties in]{On the Transfer of Completeness and Projection Properties in Truncated Vector Lattices}
\author{Mohamed Habibi, Hamza Hafsi}
\date{Source: arXiv:2505.24484v1 (CC BY 4.0)}
\begin{document}
\maketitle

\noindent\emph{Source and licence.} On the Transfer of Completeness and Projection Properties in Truncated Vector Lattices. Version arXiv:2505.24484v1 (CC BY 4.0), distributed under the Creative Commons Attribution 4.0 International licence, as recorded in the arXiv licence metadata for this exact version and shown on the abstract page https://arxiv.org/abs/2505.24484v1. Full rights notice: licenses/arxiv-2505.24484-CC-BY-4.0.txt.\par\smallskip

\noindent\textit{Editorial scope.} Counted unit: the Theorem whose source label is \texttt{None} in the article (source0.tex lines 1369--1387, source label \texttt{None}), delivered together with the article's own complete original proof of it (source0.tex lines 1389--1514, one \texttt{proof} environment of 126 lines). No mathematics is written, repaired or re-derived: statement and proof are the article's own text line for line. Required context retained uncounted: none. Every definition, display and label of those lines is retained unchanged. Omitted from this package and not redistributed: 1--1368, 1388--1388, 1515--1528. Nothing inside a retained excerpt is omitted or altered: every retained body is delivered exactly as the article's lines are written, so no omission receipt of any kind accompanies this package. Citations: the retained excerpts carry no citation command at all, so no bibliography entry of the article is transcribed and no reference is redistributed. Reference adaptation: none was needed and none was made; every reference inside the delivered unit resolves inside it, so no unresolved reference or citation is produced. Printed slips kept literal: no printed word, symbol, hypothesis or number of the counted statement or of its proof was altered. Layout and class: the article's own class is \texttt{birkau} with packages including amsmath, amssymb, url, enumitem; the wrapper is the lane's pinned \texttt{amsart} editorial wrapper at 10pt on A4, which restates the theorem-like environments the retained text needs and adds the article's own macro definitions that the retained text needs (none), with extra wrapper packages (bm, bbm, dsfont, xspace, cancel, mathtools). The delivered numbering is the unit's own. Assets: no retained span contains a figure environment, a graphics-inclusion command, a PDF-page inclusion, a table or a TikZ picture; no image file is copied into this package, the package declares no asset dependency and no image file is redistributed with it. Licence: version arXiv:2505.24484v1 of this article is distributed under the Creative Commons Attribution 4.0 International licence, as recorded in the arXiv licence metadata for this exact version and shown on the abstract page https://arxiv.org/abs/2505.24484v1. A full rights notice is delivered at licenses/arxiv-2505.24484-CC-BY-4.0.txt. Characters: every retained excerpt is pure ASCII, so the delivered engine, XeLaTeX with its default Latin Modern text font, reports no missing character.
\begin{theorem}
	Let $E$ be a truncated Riesz space.
	
	\begin{enumerate}
		\item If $E\oplus 
		%TCIMACRO{\U{211d} }%
		%BeginExpansion
		\mathbb{R}
		%EndExpansion
		$ has the projection property than $E$ has the projection property
		
		\item If $E$ is unital and has the projection property then $E\oplus 
		%TCIMACRO{\U{211d} }%
		%BeginExpansion
		\mathbb{R}
		%EndExpansion
		$ has the projection property
	\end{enumerate}
\end{theorem}

\begin{proof}
	
	\begin{enumerate}
		\item Let $B$ be a band in $E$. Then $B$ is also a band in the larger space $%
		E\oplus 
		%TCIMACRO{\U{211d} }%
		%BeginExpansion
		\mathbb{R}
		%EndExpansion
		$, and we have the orthogonal decomposition:%
		\begin{equation*}
			E\oplus 
			%TCIMACRO{\U{211d} }%
			%BeginExpansion
			\mathbb{R}
			%EndExpansion
			=B\oplus B^{d}.
		\end{equation*}
		
		Considering the ideal $B^{\prime }=B^{d}\cap E$ within $E$, we obtain the
		direct sum decomposition: 
		\begin{equation*}
			E=B\oplus B^{\prime }.
		\end{equation*}
		
		This direct sum decomposition demonstrates that $B$ is a projection band in $%
		E$.
		
		\item Let $B$ be a band in $E\oplus 
		%TCIMACRO{\U{211d} }%
		%BeginExpansion
		\mathbb{R}
		%EndExpansion
		$. Since $E$ is unital, $E$ is a band in $E\oplus 
		%TCIMACRO{\U{211d} }%
		%BeginExpansion
		\mathbb{R}
		%EndExpansion
		$, and we have the orthogonal decomposition:%
		\begin{equation*}
			E\oplus 
			%TCIMACRO{\U{211d} }%
			%BeginExpansion
			\mathbb{R}
			%EndExpansion
			=E\oplus E^{d}.
		\end{equation*}
		
		There exist two ideal $B_{1}$ and $B_{2}$ in $E$ and $E^{d}$ respectively
		such that 
		\begin{equation*}
			E=\left( B\cap E\right) \oplus B_{1}\text{ and }E^{d}=\left( B\cap
			E^{d}\right) \oplus B_{2}.
		\end{equation*}
		
		Combining these decompositions, we get:%
		\begin{eqnarray*}
			E\oplus 
			%TCIMACRO{\U{211d} }%
			%BeginExpansion
			\mathbb{R}
			%EndExpansion
			&=&\left( B\cap E\right) \oplus B_{1}\oplus \left( B\cap E^{d}\right) \oplus
			B_{2} \\
			&=&\left( \left( B\cap E\right) \oplus \left( B\cap E^{d}\right) \right)
			\oplus \left( B_{1}\oplus B_{2}\right) .
		\end{eqnarray*}
		
		Now, let's show that $B=\left( B\cap E\right) \oplus \left( B\cap
		E^{d}\right) .$ To this end, let $x\in B.$ We can decompose $x$ as $%
		x=x^{+}-x^{-}$. Further, we can decompose $x^{+}$ and $x^{-}$ based on their
		components in $E$ and $E^{d}:$%
		\begin{equation*}
			x^{+}=a_{1}+a_{2}\text{ and }x^{-}=b_{1}+b_{2}
		\end{equation*}
		where,
		
		
			\begin{itemize}
				\item $a_{1}=\sup \left( z\in E;0\leq z\leq x^{+}\right) \in E$: The
				component of $x^{+}$ in $E$
				
				\item $b_{1}=\sup \left( z\in E^{d};0\leq z\leq x^{-}\right) \in E:$ The
				component of $x^{-}$ in $E$
				
				\item $a_{2}=\sup \left( z\in E;0\leq z\leq x^{+}\right) \in E^{d}$: The
				component of $x^{+}$ in $E^{d}$
				
				\item $b_{2}=\sup \left( z\in E;0\leq z\leq x^{-}\right) \in E^{d}$: The
				component of $x^{-}$ in $E^{d}$
			\end{itemize}
	
		
		Since $x^{+}\in B$, $0\leq a_{1}\leq x^{+}$ and $B$ is an ideal, we have $%
		a_{1}\in E\cap B.$ Similarly, $b_{1}\in E\cap B.$ 
		
		Therefore,  the component of $x$ in $E$ is $a_{1}-b_{1}\in B\cap E.$
		
		Analogously,  the component of $x$ in $E^{d}$ is $a_{2}-b_{2}\in B\cap E^{d}.$
		
		Combining these components, we have:%
		\begin{equation*}
			x=a_{1}-b_{1}+a_{2}-b_{2}\in \left( B\cap E\right) \oplus \left( B\cap
			E^{d}\right) 
		\end{equation*}
		
	Hence, $B=\left( B\cap E\right) \oplus \left( B\cap E^{d}\right) $, and
		we conclude that:%
		\begin{equation*}
			E\oplus 
			%TCIMACRO{\U{211d} }%
			%BeginExpansion
			\mathbb{R}
			%EndExpansion
			=B\oplus \left( B_{1}\oplus B_{2}\right) 
		\end{equation*}
		
		This direct sum decomposition shows that $B$ is a projection band in $%
		E\oplus 
		%TCIMACRO{\U{211d} }%
		%BeginExpansion
		\mathbb{R}
		%EndExpansion
		$.
	\end{enumerate}
\end{proof}

\end{document}
